\chapter{AIと社会変革}
\section{技術能力と社会の選択}
\label{theory:ai:society}
AIが予測や生成を行えることと、その利用が社会に望ましいことは別の判断である。2017年内閣府報告書は倫理・法・経済・教育・社会・研究開発の六つの観点を置き、実例から課題を整理した\parencite{ai_society2017}。人の感情や判断への気付かない誘導、創作の価値の評価、データ利用と権利保護、利用能力の差による新たな格差は、その異なる側面である。

例えば推薦サービスは利用者に役立つ情報を選ぶ一方、見える情報の範囲を狭める場合がある。精度という指標だけでは、誰が目的を決め、何を推薦対象から除き、利用者へどのような選択の自由を与えるかは決まらない。AIを使わない選択にも配慮し、人とシステムの役割分担を具体的に設計する。

\section{AI利活用の十原則と主体ごとの対策}
\label{theory:ai:principles}
2019年総務省ガイドラインは、適正利用・適正学習・連携・安全・セキュリティ・プライバシー・尊厳自律・公平性・透明性・説明責任の十原則を整理した\parencite{ai_use2019}。適正利用は目的と役割、適正学習は入力データの質、連携はネットワーク化によるリスク、最後の四つは人の判断や待遇と説明の責任へ目を向ける。

開発者だけでなくサービス提供者、業務利用者、消費者、データ提供者で必要な対策は異なる。用途と文脈に応じて原則と措置を選ぶソフトローであり、全利用へ全措置を同じ形で要求する法律ではない。透明性も、全てのソースコードや個人データを公開する義務と同じではない。検証と説明ができる範囲を、秘密とプライバシーの保護とともに検討する。

\section{知識集約社会と対話型生成AI}
\label{theory:ai:knowledge_society}
2024年白書は、知識を価値創造の中心とし、スマート化によって製品やサービスの高付加価値化が進む社会への変化を述べる\parencite{mext_ai_whitepaper2024}。生成AIの対話型インターフェースは専門家以外にも接近しやすい入口となる。ただし入口の容易さは出力の正確さを保証しない。利用者は目的と根拠を確認し、誤りの影響が大きい判断では適切な検証を行う。

\paragraph{出題傾向の読み方}
本章の五問は、AIと社会変革について異なる資料と時点を扱う。これだけでAIが毎年必ず出題されるとはいえない。政策文書の語句と、技術が社会へ及ぼす影響を考える枠組みを分けて学ぶ。
\section{Society5.0：技術融合と人間中心の目的}
\label{theory:society:five}
狩猟社会、農耕社会、工業社会、情報社会に続く第五の社会という整理を読む。サイバー空間と現実空間の融合によって、経済的発展と社会的課題の解決の両立を目指す。中心に置くのは人間であり、技術の導入量だけを目的にしない\parencite{eth54_society2021}。

2016年の提案と2021年の計画を分ける。2021年には持続可能性、強靱性、安全・安心、多様な幸福への説明が具体化される。人文・社会科学も含む知識をつなぎ、倫理・法・社会への影響を検討する。過去問の穴埋めでは、社会段階、社会的課題、人間中心、持続可能性を別の空欄として確かめる。後年の政策用語を過去の指定語へ無条件に置き換えない。

\section{AIの社会的論点と知識集約社会}

\subsection{令和元年度再試験 問15}
\label{q:r01-ret-apt-15}

\noindent\textbf{原典：} sources/R01\_retest/aptitude.pdf、PDF物理16頁。

\textbf{公式正答：} 1（single）。

\paragraph{問題の要旨} AIの感情誘導・価値評価・データ保護・技術格差を2017年の六論点から判定する。

\Needspace{4\baselineskip}\paragraph{確認した条件・記述（独自要約）}\mbox{}\par

\noindent\textbf{原典確認：} 問題幹と各記述の主体・否定・全称・表の5候補を保存原典画像で確認。\par

\noindent\textbf{制度時点：} R01年度の原典表現を保持。現行資料の条番号や仕様を過去へ補わない。\par

\noindent\textbf{各記述の判断：} ア：○：本人が気付かない感情等の誘導は倫理課題である。；イ：○：生産性と創作価値の多様性を考える必要がある。；ウ：○：プライバシー保護と利活用の過度な萎縮防止を両立させる。；エ：○：AIリテラシーの差によるデバイドが起こり得る。\par

\noindent\textbf{前提：} 2017年報告書は倫理・法・経済・教育・社会・研究開発の六観点で便益と課題を検討する。社会の価値と機械の精度は別であり、利用できる能力の格差も課題となる。\par

\paragraph{問われる核心} AIの感情誘導・価値評価・データ保護・技術格差を2017年の六論点から判定する。

\paragraph{必要知識} 2017年報告書は倫理・法・経済・教育・社会・研究開発の六観点で便益と課題を検討する。社会の価値と機械の精度は別であり、利用できる能力の格差も課題となる。

\paragraph{考え方・正答を選ぶ理由（査読待ち初稿）} 四記述は報告書の倫理・法・社会の検討に対応する。不適切0個、1である。生産性向上は健全な利活用への期待であり、どの導入でも必ず達成されるという定理ではない。格差の可能性を、利用しない人への強制と取り違えない。

\Needspace{4\baselineskip}\paragraph{選択肢の検討（査読待ち初稿）}\mbox{}\par

\noindent\textbf{1：適切} 0。不適切なものは0個であり、一致する。\par

\noindent\textbf{2：不適切} 1。不適切なものは0個であり、この個数とは一致しない。\par

\noindent\textbf{3：不適切} 2。不適切なものは0個であり、この個数とは一致しない。\par

\noindent\textbf{4：不適切} 3。不適切なものは0個であり、この個数とは一致しない。\par

\noindent\textbf{5：不適切} 4。不適切なものは0個であり、この個数とは一致しない。\par

\paragraph{再利用する考え方} AIの技術能力、制度、価値判断、利用の機会を別軸で考える。

\paragraph{共通理論} \hyperref[theory:ai:society]{第\ref*{theory:ai:society}節}

\par 参考：\cite{ai_society2017}。

\paragraph{関連問題} \hyperref[q:r03-reg-apt-06]{R03-REG-APT-06}

\par 難易度 3（校準初稿）、優先度未評価。

\subsection{令和7年度 問7}
\label{q:r07-reg-apt-07}

\noindent\textbf{原典：} sources/R07/aptitude.pdf、PDF物理7頁。

\textbf{公式正答：} 1（single）。

\paragraph{問題の要旨} 2024年白書の知識集約社会と生成AIの普及を四つの語句へ対応させる。

\Needspace{4\baselineskip}\paragraph{確認した条件・記述（独自要約）}\mbox{}\par

\noindent\textbf{原典確認：} 問題幹と各記述の主体・否定・全称・表の5候補を保存原典画像で確認。\par

\noindent\textbf{制度時点：} R07年度の原典表現を保持。現行資料の条番号や仕様を過去へ補わない。\par

\noindent\textbf{各記述の判断：} ア：知識；イ：高付加価値化；ウ：医療；エ：対話\par

\noindent\textbf{前提：} 白書は知識を価値創造の中心とする社会、製品やサービスの高付加価値化、医療への影響、対話型インターフェースの普及を説明する。\par

\paragraph{問われる核心} 2024年白書の知識集約社会と生成AIの普及を四つの語句へ対応させる。

\paragraph{必要知識} 白書は知識を価値創造の中心とする社会、製品やサービスの高付加価値化、医療への影響、対話型インターフェースの普及を説明する。

\paragraph{考え方・正答を選ぶ理由（査読待ち初稿）} 2024年白書の対象本文は知識集約型、高付加価値化、医療、対話型生成AIと記している。四語が全て合う1である。AIの利活用の一般的な倫理課題が存在することから、本文の医療という産業分野を倫理へ置き換えない。

\Needspace{4\baselineskip}\paragraph{選択肢の検討（査読待ち初稿）}\mbox{}\par

\noindent\textbf{1：適切} 知識・高付加価値化・医療・対話。四空欄が白書の対象記述に対応する。\par

\noindent\textbf{2：不適切} 情報・標準規定化・倫理・対話。知識集約・高付加価値化・医療の三語が違う。\par

\noindent\textbf{3：不適切} 知識・標準規定化・医療・学習。高付加価値化と対話型の二語が違う。\par

\noindent\textbf{4：不適切} 情報・標準規定化・医療・対話。知識集約と高付加価値化の二語が違う。\par

\noindent\textbf{5：不適切} 知識・高付加価値化・倫理・学習。医療と対話型の二語が違う。\par

\paragraph{再利用する考え方} 政策資料の時点と対象箇所を固定し、一般的な連想で空欄を補わない。

\paragraph{共通理論} \hyperref[theory:ai:knowledge_society]{第\ref*{theory:ai:knowledge_society}節}

\par 参考：\cite{mext_ai_whitepaper2024}。

\par 難易度 3（校準初稿）、優先度未評価。

\section{AI利活用の十原則}

\subsection{令和3年度 問6}
\label{q:r03-reg-apt-06}

\noindent\textbf{原典：} sources/R03/aptitude.pdf、PDF物理6頁。

\textbf{公式正答：} 1（single）。

\paragraph{問題の要旨} 2019年AI利活用ガイドラインの十原則に含まれないものを数える。

\Needspace{4\baselineskip}\paragraph{確認した条件・記述（独自要約）}\mbox{}\par

\noindent\textbf{原典確認：} 問題幹と各記述の主体・否定・全称・表の5候補を保存原典画像で確認。\par

\noindent\textbf{制度時点：} R03年度の原典表現を保持。現行資料の条番号や仕様を過去へ補わない。\par

\noindent\textbf{各記述の判断：} ア：適正利用；イ：適正学習；ウ：連携；エ：安全；オ：セキュリティ；カ：プライバシー；キ：尊厳・自律；ク：公平性；ケ：透明性；コ：アカウンタビリティ\par

\noindent\textbf{前提：} 利用者にはサービス提供者・業務利用者・消費者が含まれる。適正利用、適正学習、連携、安全、セキュリティ、プライバシー、尊厳自律、公平性、透明性、説明責任の十原則を整理する。\par

\paragraph{問われる核心} 2019年AI利活用ガイドラインの十原則に含まれないものを数える。

\paragraph{必要知識} 利用者にはサービス提供者・業務利用者・消費者が含まれる。適正利用、適正学習、連携、安全、セキュリティ、プライバシー、尊厳自律、公平性、透明性、説明責任の十原則を整理する。

\paragraph{考え方・正答を選ぶ理由（査読待ち初稿）} 十記述は2019年ガイドライン9--10頁の十原則に全て対応する。含まれないもの0個、1である。ただしガイドライン5--6頁は利用の目的・文脈に応じた原則の選定を説明しており、全ての利用者へ全措置を同じ形で要求する法律ではない。

\Needspace{4\baselineskip}\paragraph{選択肢の検討（査読待ち初稿）}\mbox{}\par

\noindent\textbf{1：適切} 0。含まれないなものは0個であり、一致する。\par

\noindent\textbf{2：不適切} 1。含まれないなものは0個であり、この個数とは一致しない。\par

\noindent\textbf{3：不適切} 2。含まれないなものは0個であり、この個数とは一致しない。\par

\noindent\textbf{4：不適切} 3。含まれないなものは0個であり、この個数とは一致しない。\par

\noindent\textbf{5：不適切} 4。含まれないなものは0個であり、この個数とは一致しない。\par

\paragraph{再利用する考え方} 原則名の列挙と、各主体へ必要な具体策の選定を区別する。

\paragraph{共通理論} \hyperref[theory:ai:principles]{第\ref*{theory:ai:principles}節}

\par 参考：\cite{ai_use2019}。

\paragraph{関連問題} \hyperref[q:r01-ret-apt-15]{R01-RET-APT-15}

\par 難易度 3（校準初稿）、優先度未評価。

\section{Society5.0と人間中心の目的}

\subsection{令和4年度 問4}
\label{q:r04-reg-apt-04}

\noindent\textbf{原典：} sources/R04/aptitude.pdf、PDF物理3頁。

\textbf{公式正答：} 1（single）。

\paragraph{問題の要旨} Society5.0の歴史的位置と人間中心の目的。

\Needspace{4\baselineskip}\paragraph{確認した条件・記述（独自要約）}\mbox{}\par

\noindent\textbf{原典：} 問題文・否定・全五候補と組合せ表を保存原典36画像で確認。\par

\noindent\textbf{時点：} 令和4年度の本文表現を保持。後の制度改正を遡及しない。\par

\paragraph{問われる核心} Society5.0の歴史的位置と人間中心の目的。

\paragraph{必要知識} 社会の五段階、サイバー空間と現実空間、人間中心、経済と社会的課題、持続可能性。

\paragraph{考え方・正答を選ぶ理由（査読待ち初稿）} 狩猟社会1.0、農耕社会2.0、工業社会3.0、情報社会4.0に続く5.0である。技術を融合する目的は経済発展と社会的課題の解決の両立で、人間中心を核に持続可能な社会を目指す。五語の組合せを選び、2021計画の背景を後年の新計画へ置き換えない。

\Needspace{4\baselineskip}\paragraph{選択肢の検討（査読待ち初稿）}\mbox{}\par

\noindent\textbf{1：条件内で適切} 狩猟・農耕・社会・人間・持続可能。各文又は定義の判定を先に行う。全項目が一致する。\par

\noindent\textbf{2：不適切} 農耕・狩猟・社会・人間・持続可能。各文又は定義の判定を先に行う。並びが判定 狩猟・農耕・社会・人間・持続可能 と異なる。\par

\noindent\textbf{3：不適切} 狩猟・農耕・社会・人間・即応。各文又は定義の判定を先に行う。並びが判定 狩猟・農耕・社会・人間・持続可能 と異なる。\par

\noindent\textbf{4：不適切} 農耕・狩猟・技術・自然・即応。各文又は定義の判定を先に行う。並びが判定 狩猟・農耕・社会・人間・持続可能 と異なる。\par

\noindent\textbf{5：不適切} 狩猟・農耕・技術・自然・即応。各文又は定義の判定を先に行う。並びが判定 狩猟・農耕・社会・人間・持続可能 と異なる。\par

\paragraph{再利用する考え方} 社会の五段階、サイバー空間と現実空間、人間中心、経済と社会的課題、持続可能性。 設問の否定・全称・基準年を確かめ、根拠と評価範囲を分けて判断する。

\paragraph{共通理論} \hyperref[theory:society:five]{第\ref*{theory:society:five}節}

\par 参考：\cite{eth54_society2021}。

\paragraph{関連問題} \hyperref[q:r07-reg-apt-09]{R07-REG-APT-09}

\par 難易度 2（校準初稿）、優先度未評価。

\subsection{令和7年度 問9}
\label{q:r07-reg-apt-09}

\noindent\textbf{原典：} sources/R07/aptitude.pdf、PDF物理9頁。

\textbf{公式正答：} 4（single）。

\paragraph{問題の要旨} Society5.0の歴史的位置と人間中心の目的。

\Needspace{4\baselineskip}\paragraph{確認した条件・記述（独自要約）}\mbox{}\par

\noindent\textbf{原典：} 問題文・否定・全五候補と組合せ表を保存原典36画像で確認。\par

\noindent\textbf{時点：} 令和7年度の本文表現を保持。後の制度改正を遡及しない。\par

\paragraph{問われる核心} Society5.0の歴史的位置と人間中心の目的。

\paragraph{必要知識} 社会の五段階、サイバー空間と現実空間、人間中心、経済と社会的課題、持続可能性。

\paragraph{考え方・正答を選ぶ理由（査読待ち初稿）} 狩猟社会1.0、農耕社会2.0、工業社会3.0、情報社会4.0に続く5.0である。技術を融合する目的は経済発展と社会的課題の解決の両立で、人間中心を核に持続可能な社会を目指す。五語の組合せを選び、2021計画の背景を後年の新計画へ置き換えない。

\Needspace{4\baselineskip}\paragraph{選択肢の検討（査読待ち初稿）}\mbox{}\par

\noindent\textbf{1：不適切} 狩猟・農耕・社会・人間・即応。各文又は定義の判定を先に行う。並びが判定 狩猟・農耕・社会・人間・持続可能 と異なる。\par

\noindent\textbf{2：不適切} 農耕・狩猟・技術・自然・持続可能。各文又は定義の判定を先に行う。並びが判定 狩猟・農耕・社会・人間・持続可能 と異なる。\par

\noindent\textbf{3：不適切} 農耕・狩猟・社会・自然・即応。各文又は定義の判定を先に行う。並びが判定 狩猟・農耕・社会・人間・持続可能 と異なる。\par

\noindent\textbf{4：条件内で適切} 狩猟・農耕・社会・人間・持続可能。各文又は定義の判定を先に行う。全項目が一致する。\par

\noindent\textbf{5：不適切} 狩猟・農耕・技術・人間・持続可能。各文又は定義の判定を先に行う。並びが判定 狩猟・農耕・社会・人間・持続可能 と異なる。\par

\paragraph{再利用する考え方} 社会の五段階、サイバー空間と現実空間、人間中心、経済と社会的課題、持続可能性。 設問の否定・全称・基準年を確かめ、根拠と評価範囲を分けて判断する。

\paragraph{共通理論} \hyperref[theory:society:five]{第\ref*{theory:society:five}節}

\par 参考：\cite{eth54_society2021}。

\paragraph{関連問題} \hyperref[q:r04-reg-apt-04]{R04-REG-APT-04}

\par 難易度 2（校準初稿）、優先度未評価。

